\begin{center}  \Large{\textbf{ΠΕΡΙΛΗΨΗ}}  \end{center}
\begin{center}  \Large{\textbf{Βραχυπρόθεσμη Πρόβλεψη Τιμών και Φορτίου Ηλεκτρικής Ενέργειας\\
στην Ελληνική Αγορά με Μεθόδους Μηχανικής Μάθησης}}  \end{center}
\vspace{3pt}
\vspace{10pt}
\noindent
Η παρούσα διπλωματική εργασία πραγματεύεται τη βραχυπρόθεσμη πρόβλεψη τιμής
Αγοράς Επόμενης Ημέρας (DAM) και ηλεκτρικού φορτίου στην ελληνική αγορά
ηλεκτρισμού, χρησιμοποιώντας αλγόριθμους Μηχανικής Μάθησης. Κεντρικό αντικείμενο
αποτελεί η συστηματική σύγκριση τεσσάρων στρατηγικών πολυβήματης πρόβλεψης:
Teacher-Forcing (TF), Recursive/Open-Loop (OL), Multi-Input Multi-Output (MIMO)
και Walk-Forward Monthly Retraining (MR). Οι στρατηγικές εφαρμόζονται σε πέντε
αλγόριθμους (LightGBM, XGBoost, Random Forest, SVR, MLP) και αξιολογούνται σε
τρεις χρονικές περιόδους (Δεκέμβριος 2025, Q4 2025, Q1 2026) με δεδομένα από
την πλατφόρμα ENTSO-E Transparency.\\\\
\noindent
Τα πειραματικά αποτελέσματα αναδεικνύουν την Walk-Forward Monthly Retraining ως
τη βέλτιστη στρατηγική για μακροχρόνιους ορίζοντες: στο Q1 2026, το Ensemble
TF-Best3 (MR) επιτυγχάνει MAE 10,49~€/MWh για τιμή (−27,7\% vs static TF) και
67,7~MW για φορτίο (−30,0\%), αντιστεκόμενο αποτελεσματικά στη μεταβολή
κατανομής τιμών που χαρακτήρισε τον Φεβρουάριο 2026. Η ανωτερότητα αυτή
επιβεβαιώνεται στατιστικά μέσω ελέγχου Diebold-Mariano με διόρθωση
Harvey-Leybourne-Newbold. Πρόσθετα ευρήματα: τα dense intraday lags (lags 4--23)
βελτιώνουν MIMO αλλά βλάπτουν Recursive (dense lags paradox), ενώ το SVR
αποτυγχάνει σε volatile αγορές λόγω ευαισθησίας στη μεταβολή κατανομής.\\
\vspace{0.3cm}
\\ \textit{Λέξεις-κλειδιά} --- \textbf{πρόβλεψη τιμών ηλεκτρισμού, LightGBM, XGBoost, MIMO, walk-forward retraining, scheduled sampling, Diebold-Mariano}
